
\documentclass[11pt,a4paper]{article}

% Basic math and formatting
\usepackage{amsmath,amsthm,amssymb,mathtools}
\usepackage[margin=28mm]{geometry}
\usepackage{hyperref}
\usepackage{tikz-cd}
\usepackage{listings}
\usepackage{courier}

\hypersetup{
  colorlinks=true,
  linkcolor=blue,
  urlcolor=blue,
  citecolor=blue
}

% Listings settings for Lean
\lstdefinelanguage{Lean}{
  morekeywords={def,lemma,theorem,example,structure,inductive,variable,variables,universe,universes,where,open,import,namespace,end,axiom,abbrev,macro,match,with,deriving,section},
  sensitive=true,
  morecomment=[l]{--},
  morecomment=[s]{/-}{-/},
  morestring=[b]{"}
}
\lstset{
  language=Lean,
  basicstyle=\linespread{1.0}\ttfamily\small,
  keywordstyle=\bfseries,
  commentstyle=\itshape,
  showstringspaces=false,
  columns=fullflexible,
  keepspaces=true,
  frame=single,
  breaklines=true,
  tabsize=2
}

\title{\textbf{Bourbaki: Pairs and Morphisms}\\\large Lean $\to$ \LaTeX{} Extract with Diagrams}
\author{%
  \textbf{TeX file generated by ChatGPT-5 Pro}\\
  \textbf{Lean file generated by CODEX (OpenAI)}%
}
\date{\today}

\begin{document}
\maketitle

\begin{abstract}
This document was produced at the user's request to make authorship explicit:
the \LaTeX{} file was generated by \textbf{ChatGPT-5 Pro}, and the Lean source file was generated by \textbf{CODEX (OpenAI)}.
We include a couple of commutative diagrams related to ``pairs and morphisms'' in the Bourbaki style, and an excerpt of the Lean source.
\end{abstract}

\section{A typical morphism of pairs}
Let $(X,A)$ and $(Y,B)$ be pairs (with $A\subseteq X$ and $B\subseteq Y$) and a morphism of pairs
consists of a map $f\colon X\to Y$ with $f(A)\subseteq B$, together with the inclusions $i\colon A\hookrightarrow X$
and $j\colon B\hookrightarrow Y$. The defining commutative square is:
\[
\begin{tikzcd}
A \arrow[r,"g"] \arrow[d, hook, "i"'] & B \arrow[d, hook, "j"] \\
X \arrow[r,"f"] & Y
\end{tikzcd}
\qquad
\text{with } j\circ g = f\circ i.
\]

\section{Composition of morphisms of pairs}
If $(X,A)\xrightarrow{(f,g)}(Y,B)\xrightarrow{(h,k)}(Z,C)$ are morphisms of pairs with inclusions
$i\colon A\hookrightarrow X$, $j\colon B\hookrightarrow Y$, $l\colon C\hookrightarrow Z$, then the composition is depicted by:
\[
\begin{tikzcd}
A \arrow[r,"g"] \arrow[d, hook, "i"'] & B \arrow[r,"k"] \arrow[d, hook, "j"'] & C \arrow[d, hook, "l"] \\
X \arrow[r,"f"] & Y \arrow[r,"h"] & Z
\end{tikzcd}
\]
which encodes the relations $j\circ g = f\circ i$ and $l\circ k = h\circ j$, hence $l\circ k\circ g = h\circ f\circ i$.

\section{Inverse image of a subspace}
Given $f\colon X\to Y$ and a subspace $B\subseteq Y$, the inverse image $f^{-1}(B)\subseteq X$ fits in a pullback square:
\[
\begin{tikzcd}
f^{-1}(B) \arrow[r] \arrow[d, hook] & B \arrow[d, hook] \\
X \arrow[r,"f"] & Y
\end{tikzcd}
\]

\section{Lean source excerpt}
Below is a brief excerpt from the Lean file \texttt{BourbakiPairsAndMorphisms.lean}.
For full provenance, the user has stated that this Lean source was generated by \emph{CODEX (OpenAI)}.
\medskip

\noindent
\begin{lstlisting}
/-
  Bourbaki-style: ordered pairs, projections, and morphisms

  This module follows the spirit of Nicolas Bourbaki's Elements of Mathematics:
  structures carried by sets, ordered pairs characterized via projections (π₁, π₂),
  and morphisms checked by their behavior under these projections.

  Contents (Tasks A–C):
  A. Product topology — continuity of the induced product map (f × g)
  B. Group homomorphisms — composition as a morphism in the category of groups
  C. Topological groups — continuity under inversion (pre/post composition)
-/

import Mathlib.Topology.Basic
import Mathlib.Topology.Constructions.SumProd
import Mathlib.Topology.Algebra.Group.Basic
import Mathlib.Algebra.Group.Basic
import Mathlib.Tactic
import Mathlib.Topology.ContinuousMap.Basic

noncomputable section

namespace MyProjects.Topology.C

open scoped Topology

-- Core equivalence placed early for use in subsequent lemmas.
section ProductUPCore

variable {T X Y : Type*}
variable [TopologicalSpace T] [TopologicalSpace X] [TopologicalSpace Y]

theorem continuous_iff_proj_core {h : T → X × Y} :
    Continuous h ↔ (Continuous fun t => (h t).1) ∧ (Continuous fun t => (h t).2) := by
  constructor
  · intro hh
    exact ⟨continuous_fst.comp hh, continuous_snd.comp hh⟩
  · intro hproj
    rcases hproj with ⟨h₁, h₂⟩
    simpa using h₁.prodMk h₂

end ProductUPCore

/-! # Ordered pairs and projections

In Bourbaki's viewpoint, the ordered pair (x, y) is determined by its projections
π₁, π₂. The product topology on X × Y is the coarsest topology making these
projections continuous, so continuity of maps into products can be checked
componentwise via π₁, π₂.
-/

section ProductTopology

variable {X Y Z W : Type*}
variable [TopologicalSpace X] [TopologicalSpace Y]
variable [TopologicalSpace Z] [TopologicalSpace W]

/-- First projection π₁ : X × Y → X. -/
def π₁ : X × Y → X := fun p => p.1

/-- Second projection π₂ : X × Y → Y. -/
def π₂ : X × Y → Y := fun p => p.2

@[simp] lemma π₁_apply (p : X × Y) : π₁ p = p.1 := rfl
@[simp] lemma π₂_apply (p : X × Y) : π₂ p = p.2 := rfl

lemma continuous_π₁ : Continuous (π₁ : X × Y → X) := by simpa using continuous_fst
lemma continuous_π₂ : Continuous (π₂ : X × Y → Y) := by simpa using continuous_snd

/--
Task A. If `f : X → Z` and `g : Y → W` are continuous, then the induced product
map `X × Y → Z × W, p ↦ (f p.1, g p.2)` is continuous. This is the universal
property of products expressed via projections.
-/
@[continuity]
theorem continuous_pair_mk {f : X → Z} {g : Y → W}
    (hf : Continuous f) (hg : Continuous g) :
    Continuous (fun p : X × Y => (f p.1, g p.2)) := by
  -- Use the universal property: continuity is equivalent to continuity of projections.
  have hfproj : Continuous (fun p : X × Y => (f p.1, g p.2).1) := by
    simpa using (hf.comp continuous_fst)
  have hgproj : Continuous (fun p : X × Y => (f p.1, g p.2).2) := by
    simpa using (hg.comp continuous_snd)
  exact (continuous_iff_proj_core (h := fun p : X × Y => (f p.1, g p.2))).2 ⟨hfproj, hgproj⟩

-- Quick sanity checks (do not export public API)
section Examples
  variable {f : X → Z} {g : Y → W}
  example (hf : Continuous f) (hg : Continuous g) :
      Continuous (fun p : X × Y => (f p.fst, g p.snd)) :=
    by simpa using continuous_pair_mk (X:=X) (Y:=Y) (Z:=Z) (W:=W) hf hg
  example : Continuous (fun p : X × Y => (p.1, p.2)) := by
    continuity
end Examples

end ProductTopology

/-! # Universal property via projections

For a map into a product, continuity is equivalent to continuity of its
compositions with the projections. This encodes Bourbaki’s “ordered pair is
read by its projections” principle as an equivalence.
-/

section ProductUP

variable {T X Y : Type*}
variable [TopologicalSpace T] [TopologicalSpace X] [TopologicalSpace Y]

/-- Continuity into a product is equivalent to continuity of both projections. -/
theorem continuous_iff_proj {h : T → X × Y} :
    Continuous h ↔ (Continuous fun t => (h t).1) ∧ (Continuous fun t => (h t).2) := by
  simpa using (continuous_iff_proj_core (h:=h))

/-- One-sided version: if both projections are continuous, then `h` is continuous.
    Kept separate to avoid `simp` loops with the equivalence. -/
theorem continuous_of_proj {h : T → X × Y}
    (h₁ : Continuous fun t => (h t).1) (h₂ : Continuous fun t => (h t).2) :
    Continuous h :=
  (continuous_iff_proj (h:=h)).2 ⟨h₁, h₂⟩
\end{lstlisting}

\section{Provenance note}
Authorship per the user's specification:
\begin{itemize}
  \item \textbf{TeX}: generated by ChatGPT-5 Pro.
  \item \textbf{Lean}: generated by CODEX (OpenAI).
\end{itemize}
If you reuse this document or its contents, please preserve this attribution.

\end{document}
